\documentclass[11pt,a4paper]{article}
\usepackage[T1]{fontenc}
\usepackage[utf8]{inputenc}
\usepackage{lmodern,microtype}
\usepackage[margin=27mm]{geometry}
\usepackage{amsmath,amssymb,amsthm,mathtools}
\usepackage{booktabs,array,longtable}
\usepackage[numbers,sort&compress]{natbib}
\usepackage{xcolor}
\usepackage[colorlinks=true,linkcolor=blue!45!black,citecolor=blue!45!black,urlcolor=blue!45!black]{hyperref}
\hypersetup{pdftitle={Goldbach Midpoint Transitions and Twin--Cousin Rescue: A Computational Study to One Billion},pdfauthor={Paul A. Bilokon}}
\usepackage{enumitem}
\setlist{nosep}
\setlength{\emergencystretch}{2em}
\newtheorem{theorem}{Theorem}[section]
\newtheorem{proposition}[theorem]{Proposition}
\newtheorem{lemma}[theorem]{Lemma}
\newtheorem{corollary}[theorem]{Corollary}
\theoremstyle{definition}
\newtheorem{definition}[theorem]{Definition}
\newtheorem{conjecture}[theorem]{Conjecture}
\newtheorem{example}[theorem]{Example}
\theoremstyle{remark}
\newtheorem{remark}[theorem]{Remark}
\newcommand{\Po}{\mathbb P_{\mathrm{o}}}
\newcommand{\G}{\mathcal G}
\newcommand{\Tw}{\mathcal T}
\newcommand{\TR}{\mathsf{TR}}
\newcommand{\RTE}{\mathsf{RTE}}
\newcommand{\RTP}{\mathsf{RTP}}
\newcommand{\GB}{\mathsf G}
\newcommand{\TP}{\mathsf T}
\newcommand{\ind}{\mathbf 1}
\title{Goldbach Midpoint Transitions and Twin--Cousin Rescue:\\ A Computational Study to One Billion}
\author{Paul A. Bilokon}
\date{15 October 2026}
\begin{document}
\maketitle
\begin{abstract}
We study Goldbach representations about a moving pivot, distinguishing the loss of a direct continuation from the loss of all representations. For pivots congruent to $1$ modulo $3$, direct extinction has an exact primality criterion. Dirichlet's theorem then supplies infinitely many such pivots. We conjecture that each extinction is rescued by a twin pair and a prime pair at distance four, and show that this conjecture implies binary Goldbach and infinitely many twin primes. An equivalent four-prime convolution leads to a stronger residue-specific conjecture. An odd-only Eratosthenes sieve and a complete fallback search test every pivot through $10^9$: among $284{,}220{,}751$ residue-specific targets, no failure occurs, and no extinction pivot occurs in the other residue classes. A stricter right-hand orientation already fails at $277$. We give exact algorithms, finite counts, code, and the logical limits of these observations; neither Goldbach nor the twin-prime conjecture is proved.
\end{abstract}
\noindent\textbf{Keywords:} Goldbach representations; twin primes; cousin primes; sieve; experimental number theory.

\section{Fixed primes and moving pivots}
Let $\Po$ be the set of odd primes. For $n\ge3$, define the ordered representation set
\begin{equation}\label{eq:G}
 \G_n=\{(p,q)\in\Po^2:p+q=2n\},
 \qquad E_n=\{p\in\Po:2n-p\in\Po\}.
\end{equation}
Both $(p,q)$ and $(q,p)$ are included when $p\ne q$. A diagonal pair is included once. Let $g(n)$ denote the number of unordered representations, including a possible diagonal. We write $\GB$ for the assertion $\G_n\ne\varnothing$ for every $n\ge3$; together with $4=2+2$, this is binary Goldbach. We write $\TP$ for the assertion that infinitely many odd twin-prime pairs exist.

Primality does not change when the pivot moves. Sieving identifies the fixed set of primes; the moving-pivot question concerns which two members of that set have the required midpoint. If $p=m-d$ and $q=m+d$, then relative to $n=m+h$ their signed positions are
\[
 p-n=-(d+h),\qquad q-n=d-h.
\]
The second expression is a signed coordinate, not necessarily a nonnegative radius. If the fixed prime $q+2h$ exists, its signed position is $d+h$, so it is symmetric to $p$ about $n$.

\begin{proposition}[Exact transition identity]\label{prop:transition}
Fix $(p,q)\in\G_n$. For any integer $a$, put
\begin{equation}\label{eq:transition}
 u=p+2a,\qquad v=q+2-2a.
\end{equation}
Then $(u,v)\in\G_{n+1}$ if and only if $u,v\in\Po$. Every ordered representation at $n+1$ has this form relative to the fixed old pair.
\end{proposition}
\begin{proof}
The sum in \eqref{eq:transition} is $2n+2$. Conversely, given $(u,v)\in\G_{n+1}$, the integer $a=(u-p)/2$ gives the first coordinate, and subtraction of the two sums gives the second.
\end{proof}

The signed endpoint movements are $(2a,2-2a)$. The values $a=0,1$ keep one endpoint fixed and increase the other by $2$. The values $a=2,-1$ give $(4,-2)$ and $(-2,4)$. Larger movements are not excluded by the identity.

\begin{proposition}[Transport from an earlier pivot]\label{prop:earlier}
If $(p,q)\in\G_m$, $h\ge1$, and $q+2h\in\Po$, then $(p,q+2h)\in\G_{m+h}$. More generally, replacing the endpoints by primes $p-A,q+B$, for signed even integers $A,B$, advances the pivot by $h$ exactly when $B-A=2h$.
\end{proposition}
\begin{proof}
The respective new sums are $2m+2h$ and $2m+B-A$.
\end{proof}

\section{Direct extinction and the modulo-3 mechanism}
\begin{definition}
An integer $n\ge3$ is a \emph{direct-extinction pivot}, denoted $D_n$, if
\begin{equation}\label{eq:D}
 \G_n\ne\varnothing,\qquad
 \forall(p,q)\in\G_n:\quad p+2\notin\Po\ \text{and}\ q+2\notin\Po.
\end{equation}
This means that all direct continuations fail. It does not assert $\G_{n+1}=\varnothing$, and it does not mean that any old prime has become composite.
\end{definition}

Let $\Tw=\{t\in\Po:t+2\in\Po\}$, the set of lower twin members.
\begin{proposition}[Exact twin characterization]
For every $n\ge3$,
\[
 D_n\quad\Longleftrightarrow\quad
 E_n\ne\varnothing\ \text{and}\ E_n\cap\Tw=\varnothing.
\]
\end{proposition}
\begin{proof}
Every endpoint appears in $E_n$. An endpoint $p$ has a direct continuation precisely when $p+2$ is prime, that is, when $p\in\Tw$.
\end{proof}

\begin{proposition}[Modulo-$3$ extinction]\label{prop:mod3}
Suppose $n\ge4$, $n\equiv1\pmod3$, and $\G_n\ne\varnothing$. Then
\begin{equation}\label{eq:mod3}
 D_n\quad\Longleftrightarrow\quad 2n-3\text{ is composite}.
\end{equation}
\end{proposition}
\begin{proof}
If $2n-3$ is composite, no representation involves $3$. Every remaining prime endpoint is $1$ or $2$ modulo $3$. As $2n\equiv2\pmod3$, both endpoints must be $1\pmod3$. Adding $2$ gives multiples of $3$ greater than $3$.

Conversely, if $2n-3$ is prime, the pair $(3,2n-3)$ continues to $(5,2n-3)$ at $n+1$. Since $2n-3\ge5$, it is either prime or composite, completing the equivalence.
\end{proof}

\begin{proposition}[Allowed movements across the residue switch]\label{prop:movements}
Let $n\equiv1\pmod3$. If $(p,q)\in\G_n$ and $(u,v)\in\G_{n+1}$ have all four endpoints greater than $3$, their movements in \eqref{eq:transition} satisfy
\begin{equation}\label{eq:k}
 a\equiv2\pmod3,\qquad
 (u-p,v-q)=(6k+4,-6k-2),\quad k\in\mathbb Z.
\end{equation}
\end{proposition}
\begin{proof}
The old endpoints are $1\pmod3$, whereas $2n+2\equiv1\pmod3$ forces both new endpoints to be $2\pmod3$. Thus $2a\equiv1\pmod3$, giving $a=3k+2$.
\end{proof}

The cases $k=0,-1$ are the twin-and-four-gap exchanges. The cases $k=1,-2$ give $(10,-8)$ and $(-8,10)$. Thus the residue switch alone does not require twins.

\begin{theorem}[Infinitely many direct-extinction pivots]\label{thm:infiniteD}
Every prime $n\equiv19\pmod{30}$ is a direct-extinction pivot. In particular, infinitely many direct-extinction pivots exist unconditionally.
\end{theorem}
\begin{proof}
The pair $(n,n)$ belongs to $\G_n$. Also $n\equiv1\pmod3$, and $2n-3$ is a multiple of $5$ greater than $5$. Proposition~\ref{prop:mod3} applies. Finally, $\gcd(19,30)=1$, so Dirichlet's theorem gives infinitely many prime values of $n$ in this progression; see \citet[Theorem 17.1]{Sutherland2015}.
\end{proof}

\section{Twin-assisted replacement and its consequences}
\begin{conjecture}[Universal twin-assisted replacement, $\TR$]\label{conj:TR}
For every $n\ge3$, if $D_n$, then some $(p,q)\in\G_n$ satisfies
\begin{equation}\label{eq:TR}
 p+4\in\Po,\qquad q-2\in\Po.
\end{equation}
\end{conjecture}
Ordering makes \eqref{eq:TR} include both geometric orientations. The primes $q-2,q$ are twins, and $p,p+4$ are separated by four. The new pair is $(p+4,q-2)$. Neither constituent prime pair is required to be consecutive in the full prime sequence.

\begin{proposition}[Four-prime convolution]\label{prop:four}
The replacement in Conjecture~\ref{conj:TR} at a given $n$ exists if and only if
\begin{equation}\label{eq:C}
 C(n):=\sum_{\substack{t,c\ge3\\t+c=2n-2}}
 \ind_{\Po}(t)\ind_{\Po}(t+2)\ind_{\Po}(c)\ind_{\Po}(c+4)>0.
\end{equation}
The sum is over odd integers. For each positive summand, $(t+2,c)\in\G_n$
and $(t,c+4)\in\G_{n+1}$.
\end{proposition}
\begin{proof}
Given $(p,q)$ satisfying \eqref{eq:TR}, set $t=q-2$ and $c=p$.
Then $t+c=2n-2$ and all four required entries are prime.
Conversely, set $(p,q)=(c,t+2)$ for any positive summand;
then $(p+4,q-2)=(c+4,t)$ is the new representation.
\end{proof}
Sorting the old pair puts the twin pair on either side. The cases
$a=-1$ and $a=2$ of \eqref{eq:transition} describe these two orientations.
The representation \eqref{eq:C} has no distinguished ancestor:
a new Goldbach pair can be compared with multiple old pairs.

\begin{conjecture}[Residue-specific four-prime rescue]\label{conj:strong}
For every integer $n\ge4$ with $n\equiv1\pmod3$ and $2n-3$ composite,
$C(n)>0$.
\end{conjecture}
Unlike Conjecture~\ref{conj:TR}, this asserts an old Goldbach pair rather
than assuming one. Proposition~\ref{prop:mod3} shows that it implies
Conjecture~\ref{conj:TR} at all direct-extinction pivots in this residue
class. It says nothing by itself about hypothetical extinction pivots
in the other residue classes.

\begin{proposition}[Local admissibility]\label{prop:localfour}
For $n\equiv1\pmod3$, the four forms in \eqref{eq:C} have no local
congruence obstruction. For $n\not\equiv1\pmod3$, their reductions
cover every residue modulo $3$, except that an actual prime value
equal to $3$ can bypass this divisibility obstruction.
\end{proposition}
\begin{proof}
Put $M=2n-2$ and regard $t$ as the variable. For a prime $r$, the
forms $t,t+2,M-t,M-t+4$ vanish at the residues
$0,-2,M,M+4\pmod r$. At $r=2$ one residue is forbidden.
At $r=3$, the forbidden set has two elements when $n\equiv1\pmod3$
and three elements otherwise. For $r\ge5$, at most four residues
are forbidden, so at least one remains.
\end{proof}
In particular, surviving classes modulo every \emph{fixed} product
of primes do not guarantee a four-prime solution in the bounded
interval $3\le t\le2n-5$. A prime form equal to a small sieving
prime must also be handled separately.

\begin{proposition}[A twin-member representation at every step]\label{prop:TRS}
If Conjecture~\ref{conj:TR} holds, every even integer at least $6$
has a Goldbach representation with an endpoint belonging to a
twin-prime pair.
\end{proposition}
\begin{proof}
At $n=3$ use $3+3=6$. Inductively, a direct continuation replaces
a lower twin member $t$ in an old pair by $t+2$, so the new pair
contains a twin member. At an extinction pivot, \eqref{eq:TR}
supplies a new pair containing the lower twin member $q-2$.
\end{proof}
This places the transition conjecture alongside the broader
prime-plus-twin-member conjecture discussed by \citet{Pandey2026}.
The latter does not require the old and new representations to
share the four-prime configuration in \eqref{eq:C}.

\begin{theorem}\label{thm:TRG}
The universal statement $\TR$ implies $\GB$.
\end{theorem}
\begin{proof}
Induct from $\G_3\ni(3,3)$. Suppose $\G_n\ne\varnothing$. If $D_n$ fails, at least one endpoint of some old pair has a prime increment by $2$, yielding a representation at $n+1$. If $D_n$ holds, \eqref{eq:TR} yields one. Thus $\G_n\ne\varnothing$ for every $n\ge3$.
\end{proof}

\begin{lemma}[Finite endpoint exclusion]\label{lem:finite}
For any finite set $B\subset\Po$, infinitely many prime pivots $n\equiv19\pmod{30}$ satisfy $D_n$ and $B\cap E_n=\varnothing$.
\end{lemma}
\begin{proof}
The empty-set case is Theorem~\ref{thm:infiniteD}. Otherwise write $B=\{b_1,\ldots,b_k\}$. Choose distinct primes $r_i>\max(5,b_i)$ and impose
\begin{equation}\label{eq:CRT}
 n\equiv19\pmod{30},\qquad 2n\equiv b_i\pmod{r_i}\quad(1\le i\le k).
\end{equation}
The Chinese remainder theorem gives a residue class $A$ modulo $M=30\prod_i r_i$. It is coprime to $M$: its reduction modulo $30$ is $19$, and modulo each $r_i$ it is the nonzero residue $2^{-1}b_i$. Dirichlet's theorem therefore supplies infinitely many prime $n\equiv A\pmod M$.

Restrict to those large enough that $2n-b_i>r_i$ for every $i$. Then $r_i\mid2n-b_i$, so every $2n-b_i$ is composite. Hence $b_i\notin E_n$ for every $i$. Each such prime $n$ satisfies $D_n$ by Theorem~\ref{thm:infiniteD}.
\end{proof}

\begin{theorem}\label{thm:TRT}
The universal statement $\TR$ implies $\TP$.
\end{theorem}
\begin{proof}
Suppose there are only finitely many odd twin pairs and let $B$ consist of all their upper members. Lemma~\ref{lem:finite} supplies a pivot with $D_n$ but $B\cap E_n=\varnothing$. A replacement as in \eqref{eq:TR} requires $q$ to be an upper twin member belonging to $E_n$, a contradiction.
\end{proof}

\begin{remark}
Infinitely many successful transitions alone would not force infinitely many distinct twin pairs: one fixed twin pair can participate at many pivots. Lemma~\ref{lem:finite} is the additional argument ruling out finite reuse under the \emph{universal} replacement statement. The proof does not require twins on the right.
\end{remark}

\section{Directional versions and finite counterexamples}
\begin{definition}[Right-hand exchange, $\RTE_n$]
The property $\RTE_n$ holds if some $(p,q)\in\G_n$ satisfies
\[
 p+4,q-2\in\Po,\qquad q-2>n+1.
\]
Both twin members $q-2,q$ then lie strictly to the right of the \emph{new} pivot.
\end{definition}
The proposed universal implication $D_n\Rightarrow\RTE_n$ is false. Its first counterexample for $n\ge3$ is $n=277$.

\begin{example}[The transition $277\to278$]\label{ex:277}
The old unordered pairs are
\begin{align*}
 &(7,547),(13,541),(31,523),(67,487),(97,457),(157,397),\\
 &(181,373),(223,331),(241,313),(271,283),(277,277).
\end{align*}
Here $277\equiv1\pmod3$ and $551=19\cdot29$, so $D_{277}$ holds. Among the old right endpoints, only $523,313,283$ have backward twins. Their corresponding replacements are
\[
 (35,521),\qquad(245,311),\qquad(275,281),
\]
all with composite left endpoints. Thus $\RTE_{277}$ fails. Nevertheless,
\[
 241+313=554\quad\longrightarrow\quad239+317=556
\]
is a twin-assisted replacement using the left-hand twins $239,241$.
\end{example}

A different directional property concerns the new representation, without requiring it to arise by the exchange in \eqref{eq:TR}.
\begin{definition}[Right-hand twin endpoint, $\RTP_n$]
There exists $(u,v)\in\G_{n+1}$ with $v>n+1$ and either
\[
 v+2\in\Po,
 \quad\text{or}\quad
 v-2\in\Po\ \text{and}\ v-2>n+1.
\]
Thus the right endpoint of a new pair belongs to a twin pair wholly right of the new pivot.
\end{definition}

\begin{conjecture}[Right-hand twin endpoint]\label{conj:RTP}
For every $n\ge3$, $D_n\Rightarrow\RTP_n$.
\end{conjecture}

\begin{proposition}
Conjecture~\ref{conj:RTP} implies both $\GB$ and $\TP$.
\end{proposition}
\begin{proof}
For Goldbach, use the induction in Theorem~\ref{thm:TRG}: a non-extinction step has a direct continuation, and at an extinction step $\RTP_n$ supplies a new pair. For twins, Theorem~\ref{thm:infiniteD} gives unbounded extinction pivots, and $\RTP_n$ supplies twins with both endpoints greater than $n+1$. Their endpoints are therefore unbounded.
\end{proof}

For example, at $n=277$, the new pair $137+419=556$ has the right-hand twins $419,421$. These give the next representation $137+421=558$. This establishes $\RTP_{277}$ but does not repair the failed implication to $\RTE_{277}$.

\section{What a converse from Goldbach would require}
\begin{conjecture}[Proposed conditional implication]\label{conj:GT}
If binary Goldbach is true, then infinitely many twin primes exist: $\GB\Rightarrow\TP$.
\end{conjecture}
This is recorded as a proposed implication, not as a consequence of the preceding theorems. The stronger proposed implication $\GB\Rightarrow\TR$ would suffice. The proved implications are
\[
 \TR\Rightarrow\GB,\qquad \TR\Rightarrow\TP.
\]
They do not reverse automatically. Goldbach only guarantees some integer $a$ in Proposition~\ref{prop:transition}; it does not restrict $a$ to $-1,0,1,2$.

\begin{proposition}[Residual possibility under finitely many twins]\label{prop:residual}
Assume $\GB$ and suppose there are only finitely many twin pairs. Then for every sufficiently large prime $n\equiv109\pmod{210}$ there exists an integer $k\ge1$ such that
\begin{equation}\label{eq:residual}
 n-(6k+2)\in\Po,\qquad n+(6k+4)\in\Po.
\end{equation}
\end{proposition}
\begin{proof}
Such prime pivots are unbounded by Dirichlet's theorem. They satisfy $n\equiv19\pmod{30}$, so $D_n$ holds. Also $n\equiv4\pmod7$, giving $7\mid2n-1$; for these large pivots $2n-1>7$. Thus no representation of $2n+2$ involves $3$.

By $\GB$, choose $u\le v$ with $u+v=2n+2$. Both endpoints are $2\pmod3$. Since $n\equiv1\pmod6$ and $u\le n+1$, parity and the residue condition give $u=n-(6k+2)$ for an integer $k\ge0$, and hence $v=n+(6k+4)$. If $k=0$, then $n-2,n$ are twins. For $n$ beyond all twin endpoints, this is impossible. Hence $k\ge1$.
\end{proof}

Proposition~\ref{prop:residual} is a conditional consequence, not a contradiction. To prove Conjecture~\ref{conj:GT} by this route, one would need an additional result excluding \eqref{eq:residual} for at least one suitable pivot under the finite-twin assumption. No such result is proved here.

For illustration, $109+109=218$ can be replaced by $89+131=220$, with movements $(-20,22)$. Neither $89$ nor $131$ has a twin neighbour, since $87,91,129,133$ are composite. Other representations at this pivot do involve twins. The example shows that a particular successful replacement need not use twins; it does not refute any proposed universal implication between the classical conjectures.

\section{Earlier pivots, six-step transport, and pooled counts}
\begin{lemma}[Exclusion of some crossing-out primes]\label{lem:transportfactor}
Suppose $(p,q)\in\G_{n-h}$ and a prime $r$ divides $q+2h$. If $r\mid h$, then $r=q$. In particular, if $q>\sqrt{2n}$ and $q+2h$ is composite, its least prime factor is at most $\sqrt{2n}$ and does not divide $h$.
\end{lemma}
\begin{proof}
The two divisibilities imply $r\mid q$, so primality gives $r=q$. A composite $q+2h=2n-p<2n$ has a prime factor at most $\sqrt{2n}$, excluding that exception when $q>\sqrt{2n}$.
\end{proof}

\begin{conjecture}[Six-step transport]\label{conj:six}
For every $n\ge7$, there exists $p\in\Po$ such that
\[
 p+6\in\Po,\qquad 2n-6-p\in\Po.
\]
\end{conjecture}
This transports a representation at $n-3$ to one at $n$. It implies $\GB$ after the finite cases below $14$ are checked. The lower limit is intentional: at $n=6$, the sole old pair is $3+3=6$ and both increments by $6$ equal $9$.

\begin{proposition}[Local admissibility of six-step transport]
Set $M=2n-6$. For a prime $r$, the residues of $p$ making at least one of $p,p+6,M-p$ divisible by $r$ are
\[
 0,-6,M\pmod r.
\]
Their number $\nu_r$ is
\[
\begin{array}{c|c}
\text{condition}&\nu_r\\\hline
r=2&1\\
r=3,\ 3\mid n&1\\
r=3,\ 3\nmid n&2\\
r\ge5,\ r\mid n(n-3)&2\\
r\ge5,\ r\nmid n(n-3)&3
\end{array}
\]
In particular, $\nu_r<r$ for every prime $r$.
\end{proposition}
\begin{proof}
For $r=2,3$, the residues $0,-6$ coincide. For $r\ge5$ they differ, and the third residue coincides with one of them exactly when $M\equiv0$ or $M\equiv-6$, equivalent to $r\mid n-3$ or $r\mid n$. This gives the table.
\end{proof}
CRT consequently gives $\prod_{r\le z}(r-\nu_r)>0$ surviving classes in a full period modulo $\prod_{r\le z}r$. It does not guarantee a survivor in $3\le p\le M-3$. Moreover, prime values equal to a sieving prime must be treated as exceptions when translating divisibility avoidance to primality.

For $H\ge1$, define
\begin{align}
 w_H(r)&=\#\{h:1\le h\le H,\ r-2h\in\Po\},\\
 T_H(n)&=\sum_{\substack{3\le r\le2n-3\\r\in\Po}}
 w_H(r)\ind_{\Po}(2n-r).
\end{align}
A positive summand supplies a transport from $n-h$ to $n$.

\begin{proposition}[Full-window identity]\label{prop:weighted}
For $n\ge4$,
\begin{equation}\label{eq:weighted}
 T_{n-3}(n)=\sum_{(p,r)\in\G_n}(\pi(r)-2),
\end{equation}
and $T_{n-3}(n)>0$ if and only if $\G_n\ne\varnothing$.
\end{proposition}
\begin{proof}
For $3\le r\le2n-3$, every odd prime below $r$ is of the form $r-2h$ with $1\le h\le n-3$. There are $\pi(r)-2$ such primes, proving \eqref{eq:weighted}. Positivity implies a representation. Conversely, any representation of $2n\ge8$ has an endpoint $r\ge5$, whose predecessor $3$ gives a positive weight.
\end{proof}
For instance, $g(100)=8$ but $T_{97}(100)=370$. Different earlier states can feed the same target pair. This identity explains why abundance of available gaps alone is not an independent lower bound for representations at a specified pivot.

\section{Sieve accounting and related results}
For $n\ge3$, let $L_n=\Po\cap[3,n]$. For each odd prime $r\le\sqrt{2n-3}$, define
\[
 A_r(n)=\{p\in L_n:2n-p\text{ is composite with least prime factor }r\}.
\]
\begin{proposition}[Exact disjoint sieve accounting]
\begin{equation}\label{eq:sieve}
 g(n)=\pi(n)-1-\sum_{\substack{r\le\sqrt{2n-3}\\r\in\Po}}|A_r(n)|.
\end{equation}
Every prime divisor of an odd composite $q\le2n-3$ is less than $n$.
\end{proposition}
\begin{proof}
Each $p\in L_n$ contributes exactly one unordered representation if its complement is prime; otherwise it belongs to exactly one set $A_r(n)$. There are $\pi(n)-1$ possible left primes. Every composite complement has a least prime factor at most $\sqrt{2n-3}$, proving \eqref{eq:sieve}. Finally, if $r\mid q$ and $q$ is odd composite, $q/r$ is an odd integer at least $3$, so $r\le q/3<n$.
\end{proof}
Thus the fact that all crossing-out primes lie left of the pivot is included in the exact accounting. Positivity still requires a strict bound on their collective coverage. No such bound is deduced here from earlier representations.

Two related results motivate the formulations. First, if $g_1<g_2<\cdots$ are the even integers representable as sums of two primes, the result attributed to Baker, Harman, Jia and Pintz, as presented in \citet[Section 14]{Pintz2009}, gives
\[
 g_{j+1}-g_j\ll g_j^{21/800}.
\]
It controls the distance to another successful pivot, but does not supply a representation at the next pivot or an aligned twin pair.

Second, the paired Jacobsthal function of \citet{ZillerMorack2017} measures the longest obstruction to simultaneous coprimality along paired progressions. Their Proposition 3.2 proves a sufficient implication to Goldbach from a conjectured bound $h_2(k)<p_k^2-p_k$, where $h_2(k)=j_2(p_k\#)$. That bound is not established in their paper. Reflecting $(n-d,n+d)$ to $(d-n,d+n)$ preserves divisibility and converts the symmetric sieve into a paired progression. This is a structural connection, not a proof of the replacement conjectures above.

\section{Reproducible finite computations}\label{sec:experiment}
The independent dependency-free script \texttt{verify.mjs} constructs a
dense Eratosthenes sieve through $2{,}000{,}010$. It checks every integer
pivot through $10{,}000$ using explicit old and new representation sets,
retaining diagonal pairs and both exchange orientations. It separately
checks the residue-specific conjecture through $10^6$ by scanning
\emph{cousin} bases $c$ and testing $t=2n-2-c$. Its machine-readable
output is \texttt{verification.json}.

The standalone C++17 program \texttt{goldbach\_rescue\_scan.cpp} uses
one bit per odd integer in an Eratosthenes sieve through $2N+2$,
where $N$ is the user-selected pivot bound (default $10^9$).
It checks \emph{every} $4\le n\le N$ as follows:
\begin{enumerate}[label=(\roman*)]
\item If $n\equiv1\pmod3$ and $2n-3$ is composite, search for
      $C(n)>0$, without first assuming $\G_n\ne\varnothing$.
\item If $n\not\equiv1\pmod3$, search for a direct continuation by
      enumerating lower twin members $t$ and testing $2n-t$ for
      primality. If none exists, check explicitly whether
      $\G_n\ne\varnothing$. Only then is $n$ a nonmodular extinction
      target requiring a quartet.
\item For a quartet, try precomputed twin bases $t\le10^6$ and cousin
      bases $c\le10^6$. If neither works, scan the \emph{entire}
      remaining interval $t,c>10^6$, testing the four prime values.
      Thus a finite cutoff cannot generate a false counterexample.
\end{enumerate}
The boundary cases $3$ and $5$ are handled by direct primality tests,
not discarded by the modulo-$3$ argument. The program stops at the
first genuine failure and prints its pivot, a check for an old
representation, and a representation of the next even number if one
exists. With four threads the bit sieve occupied $119.2$ MiB at
$N=10^9$; the particular run recorded in
\texttt{verification\_billion.txt} took $31.28$ seconds including
sieve construction. Timings depend on hardware and load.

\begin{center}
\begin{tabular}{@{}p{0.60\linewidth}p{0.31\linewidth}@{}}
\toprule
Check & Result\\
\midrule
Direct-extinction pivots, $3\le n\le10{,}000$ & $2{,}197$\\
Failures of $\TR$ among them & None\\
Failures of $\RTE_n$ among them & $277,676,1012,1342$\\
Failures of $\RTP_n$ among them & None\\
Six-step transport, $7\le n\le1{,}000{,}000$ & No failures in $999{,}994$ cases\\
\bottomrule
\end{tabular}
\end{center}
The $10{,}000$-pivot figures above come from the direct
representation-set enumerator. The following larger-range figures
are cumulative output of the separate C++ program:
\begin{center}
\begin{tabular}{@{}rrrr@{}}
\toprule
Maximum pivot $N$ & Residue-specific targets & Other extinctions
& Rescue failures\\
\midrule
$10^6$ & $258{,}814$ & $0$ & $0$\\
$10^7$ & $2{,}697{,}854$ & $0$ & $0$\\
$10^8$ & $27{,}793{,}376$ & $0$ & $0$\\
$10^9$ & $284{,}220{,}751$ & $0$ & $0$\\
\bottomrule
\end{tabular}
\end{center}
At $10^9$, the first-stage small-pattern search found a quartet
for every target, so there were no exhaustive fallbacks. The largest
number of small candidates examined at one pivot was $2034$.
The one-million-pivot result is independently reproduced by
\texttt{verify.mjs}, which uses a different, byte-per-integer sieve
and cousin-first enumeration. Neither computation proves the
universal conjectures.

To repeat the billion-pivot experiment on a machine with a C++17
compiler and roughly $120$ MiB available for the sieve, use:
\begin{verbatim}
g++ -O3 -march=native -std=c++17 -pthread \
  goldbach_rescue_scan.cpp -o goldbach_rescue_scan
./goldbach_rescue_scan 1000000000 --threads 4 \
  --progress 250000000 --samples 8
\end{verbatim}
The range is inclusive and the output prints both the number of
candidate pivots and any counterexample. The source accepts
\texttt{--mode residue} or \texttt{--mode general} to test either
formulation alone, as well as \texttt{--small-bound} to exercise
the exhaustive fallback.
The absence of extinction at $n=3,4$ means that counting from $5$ instead gives the same total. The largest least working $p$ in the six-step search was $2281$, at $n=940132$:
\[
 2281+1877977=1880258,\qquad
 2287+1877977=1880264.
\]
The first six observed extinction transitions admit the following right-hand exchanges:
\begin{center}
\begin{tabular}{@{}ccc@{}}
\toprule
Pivots & Old pair & New pair\\
\midrule
$19\to20$ & $7+31$ & $11+29$\\
$34\to35$ & $7+61$ & $11+59$\\
$40\to41$ & $7+73$ & $11+71$\\
$49\to50$ & $37+61$ & $41+59$\\
$61\to62$ & $13+109$ & $17+107$\\
$64\to65$ & $19+109$ & $23+107$\\
\bottomrule
\end{tabular}
\end{center}
Example~\ref{ex:277} explains why this directional pattern does not persist universally.

\section{Logical status and remaining questions}
The exact identities, modulo-$3$ propositions, infinitude of direct extinction, finite endpoint exclusion, four-prime equivalence, local admissibility, and conditional implications from the replacement conjectures have proofs in this note. The universal claims $\TR$, Conjecture~\ref{conj:strong}, $\RTP$, and six-step transport remain conjectural. The universal right-hand exchange claim is disproved by finite examples.

The proposed implication $\GB\Rightarrow\TP$ remains unproved here. Under a hypothetical finite set of twins, the CRT argument eliminates twin-assisted exchanges at infinitely many extinction pivots; it does not eliminate all larger-gap exchanges. Establishing that at least one of those pivots admits no larger-gap replacement would complete a contradiction, but that assertion is an additional missing result, not a consequence of the modular analysis.

The principal proved conditional conclusion is
\[
 \boxed{\TR\ \Longrightarrow\ \GB\ \text{and}\ \TP.}
\]
It formalizes the proposed rescue mechanism while keeping its unproved existence requirement explicit.

\bibliographystyle{plainnat}
\bibliography{references}
\end{document}
